\section*{Chapitre \ref{ch:g6:earth-a-planet} --- La Terre, une planète du Système solaire}

\begin{solution}{exo:g6:earth-a-planet:1}
\underline{Mercure}, \underline{Vénus}, \underline{la Terre},
\underline{Mars}, Jupiter, Saturne, Uranus, Neptune. Les \omterm{def:g6:earth-a-planet:planet-types}{planètes
rocheuses} sont les quatre premières.
\end{solution}

\begin{solution}{exo:g6:earth-a-planet:2}
«~Une \omterm{def:g6:earth-a-planet:planet-types}{planète rocheuse} est petite et dense, avec une surface solide. Une
\omterm{def:g6:earth-a-planet:planet-types}{planète géante} est très grande, de faible masse volumique, faite surtout de
gaz et de liquides.~»
\end{solution}

\begin{solution}{exo:g6:earth-a-planet:3}
Vénus (12\,104~km, contre 12\,756~km pour la Terre). La plus grande est
Jupiter (142\,984~km)~; la plus petite est Mercure (4\,879~km).
\end{solution}

\begin{solution}{exo:g6:earth-a-planet:4}
(a) Faux~: Vénus n'a pas d'eau liquide~; il y fait bien trop chaud.
(b) Vrai ($-65$~$^\circ$C en moyenne, contre 15~$^\circ$C). (c) Faux~: les
quatre planètes les plus proches du Soleil sont les \omterm{def:g6:earth-a-planet:planet-types}{planètes rocheuses}~;
les \omterm{def:g6:earth-a-planet:planet-types}{planètes géantes} sont les quatre plus lointaines. (d) Vrai~: environ
0.7~gramme par centimètre cube.
\end{solution}

\begin{solution}{exo:g6:earth-a-planet:5}
$4879 \div 12756 \approx 0.382$, soit 38~\%~: Mercure est un peu plus
large que le tiers de la Terre.
\end{solution}

\begin{solution}{exo:g6:earth-a-planet:6}
$120536 \div 12756 \approx 9.4$~: environ 9 Terres. Saturne est une \omterm{def:g6:earth-a-planet:planet-types}{planète géante}.
\end{solution}

\begin{solution}{exo:g6:earth-a-planet:7}
La planète~A est une \omterm{def:g6:earth-a-planet:planet-types}{planète géante}~: elle est très grande et sa masse
volumique, 1.1~gramme par centimètre cube, est faible, comme celle des
\omterm{def:g6:earth-a-planet:planet-types}{planètes géantes} du Système solaire (entre 0.7 et 1.7). La planète~B est
une \omterm{def:g6:earth-a-planet:planet-types}{planète rocheuse}~: elle est petite et sa masse volumique, 5.0, est
élevée, comme celle des \omterm{def:g6:earth-a-planet:planet-types}{planètes rocheuses} (entre 3.9 et 5.5).
\end{solution}

\begin{solution}{exo:g6:earth-a-planet:8}
(a) $464 - 15 = 449$~degrés de plus. (b) $15 - (-65) = 80$~degrés de moins.
\end{solution}

\begin{solution}{exo:g6:earth-a-planet:9}
Vénus a une atmosphère très épaisse de dioxyde de carbone, qui retient la
chaleur du Soleil comme une lourde couverture et ne la laisse pas
s'échapper. Mercure, plus proche du Soleil, n'a presque pas d'air~: elle ne
peut donc pas garder sa chaleur.
\end{solution}

\begin{solution}{exo:g6:earth-a-planet:10}
$228 - 150 = 78$~millions de kilomètres de plus. $78 \div 150 = 0.52$~: la
distance de Mars est plus grande de 52~\% que celle de la Terre.
\end{solution}

\begin{solution}{exo:g6:earth-a-planet:11}
Il y a des milliards d'années, Mars avait des \omterm{def:g1:landscapes-around-us:water}{rivières}, des \omterm{def:g1:landscapes-around-us:water}{lacs} et des
\omterm{def:g5:water-shapes-the-land:flood}{crues}~: elle était plus chaude et plus humide. Aujourd'hui, son eau est
gelée, dans la glace des calottes polaires et sous le sol près des \omterm{def:g3:continents-and-oceans:poles}{pôles}.
L'eau liquide ne peut plus durer à sa surface parce que Mars est très froide
($-65$~$^\circ$C en moyenne) et parce que son air est si ténu que l'eau
gèle ou se transforme vite en vapeur.
\end{solution}

\begin{solution}{exo:g6:earth-a-planet:12}
(a) $8.9 \div 9.8 \approx 0.908$, soit 91~\%. (b) Comme une jumelle~:
presque la même taille (95~\% du diamètre de la Terre), presque la même
attraction de la pesanteur, et c'est elle aussi une \omterm{def:g6:earth-a-planet:planet-types}{planète rocheuse}. Pas
comme une jumelle~: elle est beaucoup plus chaude (464~$^\circ$C contre
15~$^\circ$C), son air est bien plus épais et fait de dioxyde de carbone
avec des \omterm{def:g4:clouds-rain-snow:cloud}{nuages} d'acide, et elle n'a pas d'eau liquide.
\end{solution}

\begin{solution}{pb:g6:earth-a-planet:1}
\textbf{1.} Vénus (108~millions de~km), la Terre (150), Mars (228).

\textbf{2.} $150 - 108 = 42$~: Vénus est 42~millions de kilomètres plus
proche. $228 - 150 = 78$~: Mars est 78~millions de kilomètres plus loin.

\textbf{3.} $6792 \div 12756 \approx 0.532$, soit 53~\%~: Mars fait environ
la moitié de la largeur de la Terre.

\textbf{4.} Toutes trois sont des \omterm{def:g6:earth-a-planet:planet-types}{planètes rocheuses}~: leurs masses
volumiques (5.2, 5.5 et 3.9 grammes par centimètre cube) sont élevées, et
leurs diamètres sont petits comparés à ceux des \omterm{def:g6:earth-a-planet:planet-types}{planètes géantes}.

\textbf{5.} L'air le plus épais est celui de Vénus (92~bars)~; le plus ténu
celui de Mars (0.01~bar).

\textbf{6.} $1 \div 0.01 = 100$~: l'air de Mars presse cent fois moins que
l'air de la Terre.

\textbf{7.} Son épais air de dioxyde de carbone retient la chaleur du Soleil
et ne la laisse pas s'échapper, comme une lourde couverture.

\textbf{8.} $92 \div 1 = 92$~: sur le sol de Vénus, l'air presse
\textbf{92 fois} plus fort que sur le sol de la Terre. Une sonde doit avoir
des parois très solides pour ne pas être écrasée (et elle doit aussi
résister à la chaleur et aux \omterm{def:g4:clouds-rain-snow:cloud}{nuages} d'acide).

\textbf{9.} La température moyenne de la Terre, 15~$^\circ$C, est comprise
entre 0~$^\circ$C et 100~$^\circ$C~: l'eau n'est ni gelée ni en
ébullition, elle reste donc liquide.

\textbf{10.} Sur Vénus, à 464~$^\circ$C, bien au-dessus de 100~$^\circ$C,
l'eau se mettrait aussitôt à bouillir et se transformerait en vapeur. Sur
Mars, à $-65$~$^\circ$C, bien en dessous de 0~$^\circ$C, elle gèlerait
(et se transformerait lentement en vapeur, parce que l'air est très ténu).

\textbf{11.} Il y a très longtemps, Mars était plus chaude et plus humide
qu'aujourd'hui~: de l'eau liquide pouvait couler à sa surface. Son \omterm{def:g5:weather-and-climate:climate}{climat} a
beaucoup changé~: c'est aujourd'hui un \omterm{def:g5:weather-and-climate:desert}{désert} froid et sec.

\textbf{12.} Par exemple~: «~Chères sondes, ici, sur la Terre, nous avons
de l'eau liquide, dans les \omterm{def:g2:water-on-earth:ocean}{océans}, les \omterm{def:g1:landscapes-around-us:water}{rivières} et les \omterm{def:g1:landscapes-around-us:water}{lacs}. Nos
températures sont douces, 15~$^\circ$C en moyenne, entre le gel et
l'ébullition de l'eau. Nous sommes à la bonne distance du Soleil, ni trop
près ni trop loin. Et notre atmosphère nous garde au chaud, donne aux êtres
vivants les gaz dont ils ont besoin et permet à l'eau de rester liquide.
Voilà pourquoi la vie est possible ici.~»
\end{solution}
